\documentclass[../../main.tex]{subfiles}
\graphicspath{{../../image/}{../../figure/}}
\begin{document}

\section{标准 Riemann 半群}\label{section:ch14-标准riemann半群}

作为上一节所建立近似解稳定性性质的推论，这里将证明：由某 $\delta$-近似解序列取 $\delta\downarrow0$ 极限所构造的系统任一解，由其初值唯一确定，并可等同于一个 $L^1$-Lipschitz 半群的轨道，该半群定义在 $L^1(-\infty,\infty)$ 的一个闭子集上。

研究的第一步是定位半群的定义域。它必须是对解正不变的集合。受\cref{section:ch14-总变差的界}分析的启发，对任一紧支集、总变差小的阶梯函数 $W(\cdot)$，我们通过如下过程关联一个数 $H(W(\cdot))$。$W(\cdot)$ 的跳跃间断通过求解经典 Riemann 问题分解为可容许激波与稀疏波的扇。在任何波碰撞发生之前，测量这些出射波的总强度 $L$ 与波相互作用势 $Q$，然后置 $H(W(\cdot))=L+2\kappa Q$，其中 $\kappa$ 是充分大正常数。设初值为 $W$ 的 $\delta$-近似解 $U$ 由\cref{section:ch14-真非线性系统的前沿追踪}的前沿追踪算法构造。按构造规则，$W$ 的所有跳跃间断将经近似 Riemann 解算器求解，因而对任一 $\delta>0$，$H(W(\cdot))$ 将与通过\eqref{align:ch14-G}定义的 Glimm 型函数 $G(t)$ 的初值 $G(0+)$ 重合。稍后当简化 Riemann 解算器起作用时，$G(t)$ 与 $H(U(\cdot,t))$ 可能彼此偏离。特别地，与 $G(t)$ 相反，$H(U(\cdot,t))$ 未必随 $t$ 非增。然而，当 $\kappa$ 充分大时，$H(U(\cdot,t))\leqslant H(U(\cdot,t-))$ 且 $H(U(\cdot,t+))\leqslant H(U(\cdot,t-))$。因而对任一 $t\geqslant0$，$H(U(\cdot,t))\leqslant H(W(\cdot))$，从而所有阶梯函数集合 $\{W(\cdot):H(W(\cdot))<r\}$ 对前沿追踪算法构造的 $\delta$-近似解是正不变的。在此准备之后，定义将充当半群定义域的集合为
\begin{align}
\mathscr{D}=\mathrm{cl}\{\text{紧支集阶梯函数 }W(\cdot):H(W(\cdot))<r\},
\label{align:ch14-domain}
\end{align}
其中 $\mathrm{cl}$ 表示在 $L^1(-\infty,\infty)$ 中的闭包。由\cref{theorem:ch01-BV-三分}，$\mathscr{D}$ 的成员是 $(-\infty,\infty)$ 上总变差以 $cr$ 为界的有界变差函数。主要结果为：



\begin{theorem}\label{theorem:ch14-标准半群}
当 $r$ 充分小时，存在一族映射 $S_t:\mathscr{D}\to\mathscr{D}$，$t\in[0,\infty)$，具有如下性质。
\begin{enumerate}[(a)]
\item $\mathscr{D}\times[0,\infty)$ 上的 $L^1$-Lipschitz 连续性：对任一 $W,\bar W\in\mathscr{D}$ 与 $t,\tau\in[0,\infty)$，
\begin{align}
\|S_tW(\cdot)-S_\tau\bar W(\cdot)\|_{L^1(-\infty,\infty)}\leqslant\kappa\{\|W(\cdot)-\bar W(\cdot)\|_{L^1(-\infty,\infty)}+|t-\tau|\}.
\label{align:ch14-semigroup-lip}
\end{align}
\item $\{S_t:t\in[0,\infty)\}$ 具有半群性质，即
\begin{align}
S_0=\text{恒等},
\label{align:ch14-S0}
\end{align}
\begin{align}
S_{t+\tau}=S_tS_\tau,\qquad t,\tau\in[0,\infty).
\label{align:ch14-semigroup}
\end{align}
\item 若 $U$ 是\eqref{align:ch13-ivp}的任一解，初值 $U_0\in\mathscr{D}$，且是某 $\delta$-近似解序列的 $\delta\downarrow0$ 极限，则
\begin{align}
U(\cdot,t)=S_tU_0(\cdot),\qquad t\in[0,\infty).
\label{align:ch14-trajectory}
\end{align}
\end{enumerate}
\end{theorem}



\begin{proof}
设 $U$ 与 $\bar U$ 是\eqref{align:ch13-ivp}的两个解，初值分别为 $U_0$ 与 $\bar U_0$，它们分别是 $\delta$-近似解序列 $\{U_{\delta_n}\}$ 与 $\{\bar U_{\bar\delta_n}\}$ 的 $\delta\downarrow0$ 极限。不假设这些近似解必然由前沿追踪算法构造。只要总变差充分小以满足\cref{section:ch14-连续依赖性}的要求，便可应用\eqref{align:ch14-lip}得
\begin{align}
\|U_{\delta_n}(\cdot,t)-\bar U_{\bar\delta_n}(\cdot,t)\|_{L^1(-\infty,\infty)}\leqslant C^2\|U_{\delta_n}(\cdot,0)-\bar U_{\bar\delta_n}(\cdot,0)\|_{L^1(-\infty,\infty)}+C\omega\max\{\delta_n,\bar\delta_n\}t.
\label{align:ch14-proof-lip}
\end{align}
令 $n\to\infty$ 取极限，我们推得
\begin{align}
\|U(\cdot,t)-\bar U(\cdot,t)\|_{L^1(-\infty,\infty)}\leqslant C^2\|U_0(\cdot)-\bar U_0(\cdot)\|_{L^1(-\infty,\infty)}.
\label{align:ch14-uniqueness-lip}
\end{align}



当 $r$ 充分小时，\cref{theorem:ch14-前沿追踪存在性}断言对任一 $U_0\in\mathscr{D}$ 可由前沿追踪算法构造的 $\delta$-近似解序列 $\{U_{\delta_n}\}$ 的极限生成\eqref{align:ch13-ivp}的解 $U$。此外，$U_\delta$ 的初值可选出使 $H(U_\delta(\cdot,0))<r$，此情形下如前所述 $H(U_\delta(\cdot,t))<r$，从而对任一 $t\in[0,\infty)$，$U(\cdot,t)\in\mathscr{D}$。由\eqref{align:ch14-uniqueness-lip}，所有这些解必重合，从而 $U$ 被唯一确定。事实上，\eqref{align:ch14-uniqueness-lip}进一步蕴含 $U$ 甚至必须与任一初值为 $U_0$、由任一 $\delta$-近似解序列取 $\delta\downarrow0$ 极限导出的解重合，无论它们是否由前沿追踪算法构造。

一旦 $U$ 如此识别，便通过\eqref{align:ch14-trajectory}定义 $S_t$。Lipschitz 连续性性质\eqref{align:ch14-semigroup-lip}由结合\eqref{align:ch14-uniqueness-lip}与\eqref{align:ch13-ivp}得到，而\eqref{align:ch14-S0}是显然的。为验证\eqref{align:ch14-semigroup}，只需注意对任一固定 $\tau>0$，$U(\cdot,\tau+\cdot)$ 是\eqref{align:ch13-ivp}以 $U(\cdot,\tau)$ 为初值的解，它由 $\delta$-近似解取 $\delta\downarrow0$ 极限导出，因而由唯一性必与 $S_tU(\cdot,\tau)$ 重合。证明完成。
\end{proof}



``标准 Riemann 半群''这一术语通常用于 $S_t$，以提醒其建筑砌块是 Riemann 问题的解。此半群是否也涵盖由其他方法导出的解，将在下一节处理。

\end{document}
